\begin{center}
\vspace*{3cm}
\textsc{\huge \bfseries Abstract}\\[1cm]
\end{center}

\begin{spacing}{1.3}

This report presents a comprehensive, end-to-end MLOps pipeline for melanoma detection on edge devices. The pipeline spans the full machine learning lifecycle -- from data ingestion through CI/CD -- applying modern software engineering practices to produce a reproducible, production-ready system.

The pipeline uses the HAM10000 dataset (10,015 dermoscopic images, 7 classes) with patient-aware GroupShuffleSplit on lesion identifier to prevent data leakage. Data quality is enforced through Great Expectations with 7 validation expectations and 100\% image integrity verification. Three convolutional architectures were trained and compared: EfficientNet-B0 achieved the highest ROC-AUC of 0.91 (4.3M parameters), followed by ResNet50 at 0.84 (24M) and MobileNetV3-Small at 0.84 (1.5M). At the clinically-calibrated threshold of 0.15, sensitivity reaches 80.6\% with specificity of 82.8\%. The best-performing model was exported to ONNX (FP32, 16.6~MB) achieving a mean inference latency of 6.7~ms and p99 latency of 10.3~ms on CPU. Post-training INT8 quantisation was investigated across five strategies but proved incompatible with EfficientNet-B0, producing 10--15\% AUC degradation due to SiLU activation and depthwise separable convolution incompatibility with integer arithmetic.

The deployment architecture centres on a FastAPI inference server with 6 endpoints, containerised through Docker multi-stage builds. CI/CD is orchestrated via GitHub Actions with 41 passing tests across 5 pipeline stages. An interactive Gradio demonstration is hosted on Hugging Face Spaces. Explainability is provided through Grad-CAM visualisation, MC-Dropout uncertainty quantification, and temperature scaling calibration. This project demonstrates that MLOps provides the engineering framework to transform a research model into a robust, edge-deployable clinical decision support tool.

\end{spacing}
